\documentclass[11pt,a4paper]{article}
\usepackage[utf8]{inputenc}
\usepackage[T1]{fontenc}
\usepackage[margin=0.72in]{geometry}
\usepackage{enumitem}
\usepackage{titlesec}
\usepackage{xcolor}
\usepackage{hyperref}
\usepackage{array}

\definecolor{primarycolor}{RGB}{33,37,41}
\definecolor{secondarycolor}{RGB}{0,102,204}
\definecolor{accentcolor}{RGB}{90,98,104}
\hypersetup{colorlinks=true,urlcolor=secondarycolor}
\titleformat{\section}{\large\bfseries\color{primarycolor}}{}{0em}{}[\titlerule]
\titlespacing*{\section}{0pt}{1.15em}{0.45em}
\setlist[itemize]{leftmargin=1.4em,itemsep=0.22em,topsep=0.25em,parsep=0pt}
\pagestyle{empty}
\setlength{\parindent}{0pt}
\setlength{\parskip}{0.25em}

\begin{document}
\begin{center}
{\Huge\bfseries\color{primarycolor} MARTA KOWAL-BRANDT}\\[0.3em]
{\large DevOps / Platform Engineering | QA Automation Engineer}\\[0.25em]
Wroclaw, Poland | \href{mailto:marta.kowalbrandt@example.com}{marta.kowalbrandt@example.com} | +48 555 303 303 | \href{https://gitlab.com/mkb-example}{gitlab.com/mkb-example}
\end{center}

\section*{PROFESSIONAL SUMMARY}
QA automation engineer with 7 years of experience, including the last two focused primarily on CI/CD pipelines, test environments and build tooling. Owns GitLab CI pipelines for three product teams and has built merge-request ephemeral environments on a shared Kubernetes cluster. Experienced in improving pipeline performance, infrastructure automation, observability for test metrics, and cross-team technical documentation. Building further infrastructure-as-code and Kubernetes experience through a personal k3s and Terraform project.

\section*{TECHNICAL SKILLS}
CI/CD and build tooling: GitLab CI, Jenkins, pipeline optimization, artifact publishing\\
Containers and orchestration: Docker, Docker Compose, Helm, Kubernetes (namespace-level access; basic), k3s\\
Infrastructure as code and automation: Terraform (personal project), Ansible (basic), Bash, Python\\
Observability and data: Prometheus, Grafana, Loki, PostgreSQL, SQL\\
Testing: pytest, Playwright, Selenium, Postman\\
Operating systems and version control: Linux, Git\\
Cloud learning: AWS EC2 and S3 basics (personal free tier)

\section*{PROFESSIONAL EXPERIENCE}
\textbf{QA Automation Engineer (Test Infrastructure) | Odrava Health Software | Wroclaw} \hfill \textit{2021–Present}
\begin{itemize}
\item Own GitLab CI pipelines for three product teams (approximately 25 developers), covering build, unit, integration and end-to-end test stages, and artifact publishing.
\item Reduced average pipeline time from 38 to 14 minutes by parallelizing test jobs, caching dependencies and separating end-to-end tests into their own stage; monthly CI runner costs fell by about 30\%.
\item Built merge-request ephemeral environments comprising an application, PostgreSQL and seeded data, deployed via Helm to the shared Kubernetes cluster from GitLab CI; environments are automatically removed after merge or 48 hours of inactivity.
\item Maintain a pytest and Playwright end-to-end suite of approximately 900 tests; reduced the flaky-test rate from 7\% to below 1\%.
\item Built a Python and PostgreSQL pipeline-results workflow and Grafana dashboards to show flaky tests by team.
\item Wrote Bash and Python tooling to seed test data in PostgreSQL.
\item Wrote the internal guide to CI pipelines used in team onboarding; mentored two QA interns in 2023.
\end{itemize}

\textbf{QA Engineer | Brightmill Software House | Wroclaw} \hfill \textit{2019–2021}
\begin{itemize}
\item Tested web and mobile applications for client projects using manual and automated testing.
\item Introduced a Selenium and Java automated regression suite for the largest client project.
\item Configured Jenkins jobs to run nightly regression tests.
\end{itemize}

\textbf{Customer Support Specialist (2nd line) | Brightmill Software House} \hfill \textit{2018–2019}
\begin{itemize}
\item Investigated customer issues in logs and databases and reproduced bugs for developers.
\end{itemize}

\section*{SELECTED PROJECTS}
\textbf{Homelab k3s Cluster with GitOps | Personal} \hfill \textit{2024–2025}
\begin{itemize}
\item Built a three-node Raspberry Pi 5 k3s cluster; deploys applications from a Git repository using Argo CD.
\item Set up Prometheus, Grafana and Loki for monitoring, with Traefik ingress and cert-manager for TLS.
\item Provisioned infrastructure for a small Hetzner VPS using Terraform; used Ansible at a basic level.
\end{itemize}

\textbf{Merge-Request Ephemeral Environments | Odrava Health} \hfill \textit{2023–2024}
\begin{itemize}
\item Delivered short-lived application and PostgreSQL environments with seeded data for each merge request, deployed with Helm to the shared Kubernetes cluster from GitLab CI.
\item Set environments to be destroyed automatically after merge or 48 hours of inactivity, enabling QA and product review before merge and end-to-end testing against the environment.
\end{itemize}

\textbf{Flaky Test Dashboard | Odrava Health} \hfill \textit{2023}
\begin{itemize}
\item Collected CI test results in PostgreSQL and built Grafana dashboards showing flaky tests by team.
\item The flaky-test rate fell from 7\% to below 1\% after teams began addressing the leading offenders.
\end{itemize}

\textbf{CI Pipeline Speed-up | Odrava Health} \hfill \textit{2022}
\begin{itemize}
\item Profiled pipelines across three teams, parallelized test jobs, added dependency caching and separated end-to-end testing into its own stage.
\item Reduced average pipeline time from 38 to 14 minutes and monthly CI runner costs by about 30\%.
\end{itemize}

\section*{EDUCATION}
BEng, Computer Engineering | Silesian Lowlands University of Technology | 2018

\section*{CERTIFICATION AND PROFESSIONAL DEVELOPMENT}
ISTQB Foundation Level | 2019\\
Currently studying for the Certified Kubernetes Administrator (CKA) exam; exam scheduled for January 2027.

\section*{LANGUAGES}
Polish (native), English (C1), German (A2)
\end{document}